\documentclass[11pt]{article}

\usepackage[margin=1in]{geometry}
\usepackage{amsmath, amssymb, amsthm}
\usepackage{array, booktabs}
\usepackage{enumitem}
\usepackage{hyperref}
\usepackage{setspace}
\usepackage{graphicx}
\usepackage{xcolor}
\usepackage{float}
\usepackage{caption}

\setstretch{1.15}
\graphicspath{{figures/}}

\theoremstyle{definition}
\newtheorem{definition}{Definition}

% --------------------
% Course info
% --------------------
\newcommand{\course}{ST5230: Applied Natural Language Processing}
\newcommand{\semester}{Semester 2, AY2025/2026}

\title{\textbf{Do LLM Rankings Change if You Ask Differently?}\\
\large Evaluating Ranking Reversals under Non-semantic Prompt Perturbations and Subset Resampling}
\author{Group 14\\
ZHANG QIANCHI, PENG YANGYUNZHI, JIANG YIFAN, MENG XIANGCHEN}
\date{}

\begin{document}
\maketitle
\vspace{-1em}

\begin{abstract}
Leaderboards serve as the primary metric for Large Language Model (LLM) evaluation, yet they often overlook ranking stability under minor protocol variations. This paper systematically investigates this fragility by analyzing ranking reversals under non-semantic prompt perturbations and evaluation-subset resampling. Evaluating four representative models---GPT-4o-mini, Gemini-2.0-flash, Qwen-2.5-7B, and Llama-3.1-8B---across MMLU, ARC, and HellaSwag, we introduce two diagnostic metrics: Prompt-wise Ranking Reversal Rate (PRRR) and Subset Ranking Reversal Rate (SRRR). Based on 14,400 responses and 1,000 resampling iterations, our results reveal that ranking stability is highly task-dependent: while ARC is robust, HellaSwag exhibits significant fragility with a maximum PRRR of 56.00\%. Furthermore, subject-level analysis on MMLU uncovers acute instability in fields like sociology and history, where PRRR and SRRR peak at 74.30\% and 67.90\%, respectively. We argue that current evaluations often conflate average accuracy with ranking robustness, and we advocate for stability diagnostics to distinguish true model superiority from protocol-dependent artifacts.
\end{abstract}

\section{Introduction}

Leaderboards are a convenient way to evaluate large language models (LLMs). A model is assigned a benchmark score, models are ranked based on that score, and the resulting order is often taken as a stable indication of relative capability. This practice is useful, but it also relies on a strong assumption: if the evaluation protocol is changed in a reasonable but non-semantic way, the ranking should remain largely the same.

In practice, however, this assumption does not always hold. Even when the underlying task is fixed, differences in prompt wording, instruction style, and answer format can affect model behavior. Likewise, the observed ranking may depend on which subset of examples is used for evaluation. When the performance gap between two models is small, prompt variation or subset noise may be enough to change the apparent winner.

In this project, we study LLM evaluation from the perspective of \textbf{ranking stability}. Rather than asking only which model achieves the highest average accuracy, we ask whether that ranking remains the same under prompt perturbations and subset resampling. This perspective is important because, in practice, users often care about rank comparisons: which model ranks first, which ranks second, and whether a reported improvement reflects a robust advantage or merely a point-estimate artifact.

Our contributions are as follows:
\begin{itemize}[leftmargin=*]
    \item We construct a controlled benchmark experiment with four task-preserving prompt templates across MMLU, ARC-Challenge, and HellaSwag.
    \item We distinguish between two sources of instability: prompt-induced ranking changes and subset-induced ranking changes.
    \item We report PRRR and SRRR as ranking stability diagnostics, showing that HellaSwag is much less stable than ARC-Challenge and MMLU at the dataset level.
    \item We include an auxiliary expanded MMLU subject extension and show that ranking instability at the subdomain level can be much greater than at the aggregate benchmark level.
\end{itemize}

\section{Related Work}

This work is related to three lines of prior research. First, benchmark construction studies such as MMLU, ARC, and HellaSwag have established widely used multiple-choice evaluation settings for knowledge, reasoning, and commonsense completion tasks \cite{hendrycks2020mmlu, clark2018arc, zellers2019hellaswag}. These benchmarks provide convenient and automatically scorable testbeds, but the benchmark papers do not necessarily show that model rankings remain invariant under prompt and subset variation.

Second, work on prompt robustness has shown that LLM outputs can be sensitive to prompt form even when the task semantics remain unchanged. PromptRobust, for instance, examines robustness under adversarial prompt perturbations and shows that prompt design can meaningfully affect downstream performance \cite{zhu2023promptrobust}. Our work differs in that we focus on whether such variation changes the \emph{ordering} between models, rather than only their individual scores.

Third, recent discussions on evaluation reliability have emphasized uncertainty, repeated trials, and error bars in LLM evaluation \cite{miller2024errorbars, fourrier2024openllm}. We build on this perspective by treating the leaderboard itself as an empirical object whose robustness should be examined. In this sense, our project is closer to a stress test of leaderboard interpretation than a search for the single best-performing model.

\section{Experimental Setup}

\subsection{Research Question}

We study the following testable question:
\begin{quote}
    \textbf{Does model ranking remain stable under task-preserving prompt perturbations and evaluation-subset variation?}
\end{quote}

We operationalize this question using two types of ranking reversal. Prompt-wise ranking reversal measures whether two prompt templates induce different rankings on the same resampled item subset. Subset ranking reversal measures whether a resampled subset induces a different ranking from the full item set under the same prompt.

\subsection{Models}

We evaluate four models through the OpenRouter API:
\begin{enumerate}[leftmargin=3.75em]
    \item \texttt{openai/gpt-4o-mini} (GPT-4o-mini)
    \item \texttt{google/gemini-2.0-flash-001} (Gemini-2.0-flash)
    \item \texttt{qwen/qwen-2.5-7b-instruct} (Qwen-2.5-7B)
    \item \texttt{meta-llama/llama-3.1-8b-instruct} (Llama-3.1-8B)
\end{enumerate}

All model calls use temperature $0$ and \texttt{max\_tokens=150}. This deterministic decoding setting reduces sampling randomness so that the analysis focuses on prompt and item-subset variation.

\subsection{Datasets and Tasks}

\begin{center}
\begin{tabular}{@{}
>{\raggedright\arraybackslash}p{0.22\linewidth}
>{\raggedright\arraybackslash}p{0.12\linewidth}
>{\raggedright\arraybackslash}p{0.55\linewidth}
@{}}
\toprule
\textbf{Dataset} & \textbf{Size} & \textbf{Details} \\
\midrule
MMLU & 300 & Main dataset-level sample; 50 questions from each of 6 subjects in the test split. \\
ARC-Challenge & 300 & Science question answering benchmark; sampled from the official test split. \\
HellaSwag & 300 & Commonsense sentence-completion benchmark; sampled from validation because test labels are not public. \\
Auxiliary expanded MMLU extension & 5 $\times$ 200 & A Separate 200 items from each of the following subjects: \texttt{high\_school\_mathematics}, \texttt{conceptual\_physics}, \texttt{high\_school\_world\_history}, \texttt{philosophy}, and \texttt{sociology}. \\
\bottomrule
\end{tabular}
\end{center}

\subsection{Experimental Design}

The project has two empirical layers. The main standardized pipeline uses the repository's default Phase 1 to Phase 4 workflow: sample 900 benchmark items, expand them into four prompt templates, collect fixed model responses, parse them into binary correctness labels, and run local resampling analysis on the scored outputs. This main experiment is the source of the dataset-level results in this report.

We then add an auxiliary expanded MMLU subject extension to stress-test whether benchmark-level conclusions survive at a narrower subdomain level. This extension uses a separate subject-preparation script and a separate response file, but keeps the same four models, the same four prompt templates, and the same deterministic decoding settings as the main experiment.

\subsection{Prompt Design}

Each question is converted into four prompt templates that preserve the answer options and task semantics while varying surface form:
\begin{itemize}[label={}]
    \item \textbf{T1 Minimal}: the question and choices followed by ``Select one option.''
    \item \textbf{T2 Benchmark}: explicit multiple-choice framing, with ``Question'' and ``Options'' fields.
    \item \textbf{T3 Natural}: conversational framing, e.g., ``I have a multiple-choice question for you.''
    \item \textbf{T4 Order/Phrasing}: same task and options with altered instruction order and phrasing.
\end{itemize}

The main dataset-level experiment therefore expands 900 items into $900 \times 4 = 3600$ prompt requests. Running four models over these requests yields $3600 \times 4 = 14400$ model responses. The auxiliary expanded MMLU extension adds 1000 subject-level questions, 4000 prompt requests, and 16000 responses.

\subsection{Scoring}

All tasks are multiple-choice. We parse each model output into one of \texttt{A/B/C/D} and compare it with the gold label to obtain a binary score. In the main dataset-level experiment, there are 3 parse failures out of 14,400 responses. In the auxiliary expanded MMLU extension, there are no parse failures.

\section{Metrics and Statistical Protocol}

\subsection{Definition of Metrics}

\begin{definition}[Accuracy and Ranking]
    For any item set $A$ and template $T$, a model $M$'s accuracy on $A$ under $T$ is
    \[
    \mathrm{Acc}(A,T,M) = \frac{1}{|A|}\sum_{x\in A} s(x,T,M),
    \]
    where $s(x,T,M)\in\{0,1\}$ is the fixed binary correctness score for $M$ on item $x$ under template $T$. Further, all models (four models in our experiment) can be sorted by their accuracy on $A$ under $T$, inducing a ranking, denoted $R_{A,T}$.
\end{definition}

\begin{definition}[Prompt-wise Ranking Reversal Rate, PRRR]
    For an item set $P$, two templates $T_1$ and $T_2$, draw $N$ subsets $P_i \subset P$; the PRRR on $P$ between $T_1$ and $T_2$ is
    %$$ \text{PRRR}(P, T_1, T_2):= \frac{1}{N} \# \{i \in [1:N] : R_{P_i, T_1} \neq R_{P_i, T_2}\}.$$
    $$\text{PRRR}(P, T_1, T_2):= \frac{1}{N} \sum_{i=1}^{N} \mathbf{1}\big\{ R_{P_i,T_1}\neq R_{P_i,T_2}\big\}$$
    that is, the frequency of different rankings occuring under two templates among the $N$ subsets.
\end{definition}

\begin{definition}[Subset Ranking Reversal Rate, SRRR]
    For an item set $P$ and template $T$, draw $N$ subsets $P_i \subset P$; the SRRR on $P$ under $T$ is
    %$$\text{SRRR}(P, T) := \frac{1}{N} \# \{i \in [1:N] : R_{P,T} \neq R_{P_i,T}\}.$$
    $$\text{SRRR}(P, T) := \frac{1}{N} \sum_{i=1}^{N} \mathbf{1}\big\{R_{P,T}\neq R_{P_i,T}\big\}.$$
    that is, the frequency of different rankings occuring under $T$ when asked subsets and the full set.
\end{definition}

\iffalse
For an item set $P$, prompt template $T$, and model $m$, let

 The ranking $R_{P,T}$ is the ordering of the four models by $\mathrm{Acc}(P,T,m)$.

\subsection{Prompt-wise Ranking Reversal Rate}

Prompt-wise Ranking Reversal Rate (PRRR) captures whether two prompt templates induce different rankings on the same resampled subset. For an item set $P$, baseline template $T_1$, comparison template $T_2$, and resampled subsets $P_i\subset P$ for $i=1,\ldots,N$,
\[
\mathrm{PRRR}(P,T_1,T_2)
=
\frac{1}{N}
\sum_{i=1}^{N}
\mathbf{1}\left[R_{P_i,T_1}\neq R_{P_i,T_2}\right].
\]
In this project, $T_1$ is the minimal instruction template and $T_2$ ranges over the three non-minimal templates.

\subsection{Subset Ranking Reversal Rate}

Subset Ranking Reversal Rate (SRRR) captures whether a subset ranking differs from the full-set ranking under the same prompt template:
\[
\mathrm{SRRR}(P,T)
=
\frac{1}{N}
\sum_{i=1}^{N}
\mathbf{1}\left[R_{P,T}\neq R_{P_i,T}\right].
\]
\fi

The two types of ranking reversal rates defined above isolates two possible sources of instability. PRRR is the metric for measuring prompt-to-prompt disagreement, while SRRR is for describing the instability caused by using an incomplete evaluation subset.

\subsection{Resampling Protocol}

For each dataset or subject scope, we run $N=1000$ resampling iterations. Each iteration samples a half-sized subset without replacement from the current scope. For the main 300-item datasets, each subset contains 150 items. For each subject in the auxiliary expanded MMLU extension, each subset contains 100 of the 200 subject items. We also compute rank-probability tables as supplementary summaries, but PRRR and SRRR remain our main stability diagnostics in the report.

\section{Results}\label{section5Results}

\subsection{Dataset-level Accuracy}

Table~\ref{tab:acc} reports prompt-averaged accuracy and rank for the main dataset-level experiment. Gemini is first on ARC-Challenge and MMLU, while GPT-4o-mini is first on HellaSwag.

\begin{table}[H]
\centering
\begin{tabular}{@{}llcc@{}}
\toprule
\textbf{Dataset} & \textbf{Model} & \textbf{Accuracy} & \textbf{Rank} \\
\midrule
ARC-Challenge & Gemini-2.0-flash & 94.00\% & 1 \\
ARC-Challenge & GPT-4o-mini & 90.83\% & 2 \\
ARC-Challenge & Qwen-2.5-7B & 85.08\% & 3 \\
ARC-Challenge & Llama-3.1-8B & 77.33\% & 4 \\
\midrule
HellaSwag & GPT-4o-mini & 89.17\% & 1 \\
HellaSwag & Gemini-2.0-flash & 87.25\% & 2 \\
HellaSwag & Qwen-2.5-7B & 84.25\% & 3 \\
HellaSwag & Llama-3.1-8B & 68.33\% & 4 \\
\midrule
MMLU & Gemini-2.0-flash & 85.08\% & 1 \\
MMLU & GPT-4o-mini & 73.42\% & 2 \\
MMLU & Qwen-2.5-7B & 70.33\% & 3 \\
MMLU & Llama-3.1-8B & 60.08\% & 4 \\
\bottomrule
\end{tabular}
\caption{Prompt-averaged dataset-level accuracy.}
\label{tab:acc}
\end{table}

\begin{figure}[H]
\centering
\includegraphics[width=\linewidth]{accuracy_boxplots.png}
\caption{Subset accuracy distributions across datasets, models, and prompt templates. ARC-Challenge has relatively separated model distributions, while HellaSwag and parts of MMLU show closer margins.}
\label{fig:accuracy}
\end{figure}

\subsection{Dataset-level Ranking Reversal}

\begin{figure}[H]
\centering
\includegraphics[width=0.53\linewidth]{prompt_ranking_reversal_rate.png}
\caption{Prompt-wise ranking reversal rate (PRRR). HellaSwag has the largest prompt-induced reversal rate, especially between the minimal and benchmark-style prompts.}
\label{fig:prrr}
\end{figure}

\begin{figure}[H]
\centering
\includegraphics[width=0.53\linewidth]{subset_ranking_reversal_rate.png}
\caption{Subset ranking reversal rate (SRRR). HellaSwag is also most sensitive to using only a half-sized evaluation subset.}
\label{fig:srrr}
\end{figure}
Table~\ref{tab:dataset_rrr} summarizes the maximum PRRR and maximum SRRR observed for each benchmark. HellaSwag is the most fragile dataset in both prompt-induced and subset-induced ranking instability.

\begin{table}[H]
\small
\centering
\begin{tabular}{@{}lcc@{}}
\toprule
\textbf{Dataset} & \textbf{Max PRRR} & \textbf{Max SRRR} \\
\midrule
ARC-Challenge & 10.60\% & 11.20\% \\
HellaSwag & 56.00\% & 42.80\% \\
MMLU & 15.60\% & 17.50\% \\
\bottomrule
\end{tabular}
\caption{Maximum dataset-level PRRR and SRRR.}
\label{tab:dataset_rrr}
\end{table}



\subsection{Expanded MMLU Subject-level Analysis}

The auxiliary expanded MMLU subject extension reveals substantially stronger instability than the aggregate MMLU result.

\begin{figure}[H]
\centering
\includegraphics[width=0.9\linewidth]{mmlu_subject_expanded_accuracy_boxplots.png}
\caption{Expanded MMLU subject-level subset accuracy distributions. The subject-level view exposes more heterogeneous margins than the aggregate benchmark view.}
\label{fig:mmlu_acc}
\end{figure}

Table~\ref{tab:mmlu_subject} reports the winning model, top accuracy, maximum PRRR and SRRR for each subject.

\begin{table}[H]
\small
\centering
\begin{tabular}{@{}lcccc@{}}
\toprule
\textbf{Subject} & \textbf{Winner} & \textbf{Top Acc.} & \textbf{Max PRRR} & \textbf{Max SRRR} \\
\midrule
High-school mathematics & Gemini-2.0-flash & 62.75\% & 32.20\% & 38.70\% \\
Conceptual physics & Gemini-2.0-flash & 90.50\% & 25.60\% & 26.20\% \\
High-school world history & Gemini-2.0-flash & 91.38\% & 74.30\% & 55.00\% \\
Philosophy & Gemini-2.0-flash & 87.38\% & 56.10\% & 49.00\% \\
Sociology & GPT-4o-mini & 89.38\% & 54.30\% & 67.90\% \\
\bottomrule
\end{tabular}
\caption{Expanded MMLU subject-level stability results.}
\label{tab:mmlu_subject}
\end{table}


\begin{figure}[H]
\centering
\includegraphics[width=0.78\linewidth]{mmlu_subject_expanded_prrr.png}
\caption{Expanded MMLU subject-level PRRR. High-school world history and sociology exhibit particularly large prompt-induced ranking disagreement.}
\label{fig:mmlu_prrr}
\end{figure}

\begin{figure}[H]
\centering
\includegraphics[width=0.78\linewidth]{mmlu_subject_expanded_srrr.png}
\caption{Expanded MMLU subject-level SRRR. Sociology has the largest subset-induced instability under the minimal prompt.}
\label{fig:mmlu_srrr}
\end{figure}

\section{Analysis and Discussion}

As shown by results in section \ref{section5Results}, it is possible for variation in prompt choice and task sets to overturn performance ranking of the models. Further, different benchmarks have shown different degree of ranking stability.

\subsection{ARC-Challenge: a Robust Leaderboard Case}

ARC-Challenge behaves like an idealized, stable leaderboard interpretation the most. Gemini leads the benchmark by prompt-averaged accuracy most of the time, model margins are relatively wide, and both PRRR and SRRR remain low. The maximum PRRR and SRRR is 10.60\% and 11.20\%, respectively. This suggests that on ARC-Challenge, prompt and subset variations are usually too small to overturn the ranking.

\subsection{HellaSwag: a Fragile Leaderboard Case}

HellaSwag provides the clearest counterexample. The performance of Gemini and Qwen are close enough that rank order becomes sensitive to prompt and subset choice. HellaSwag reaches a high maximum PRRR of 56.00\% and a maximum SRRR of 42.80\%. This means that a single leaderboard constructed from one prompt and one subset may overstate the certainty of the ranks.

\subsection{MMLU: Stable Aggregate Winner, Unstable Slices}

At the master dataset level, MMLU is more stable than HellaSwag but less stable than ARC-Challenge. Gemini is clearly first in the aggregate, but the middle ranks are closer. The auxiliary subject-level extension shows why aggregation can hide instability: high-school world history reaches 74.30\% PRRR, and sociology reaches 67.90\% SRRR, while some other subjects obtained less ranking reversals, such as conceptual physics. Thus, even if the aggregate benchmark has a stable winner, subdomain-level rankings may remain fragile.

\subsection{A Concrete Reversal Case}

Sociology provides a clear example of why prompt choice matters when ranking language models. The subject-level average winner is GPT-4o-mini, by a very small margin (89.38\% vs. 89.25\%), over Gemini, while the prompt-level leader flips repeatedly. Under T1 Minimal and T3 Natural, Gemini is first; under T2 Benchmark and T4 Order/Phrasing, GPT-4o-mini is first. The instability is therefore not driven by some collapse in accuracy under some certain prompts, but by extremely small margins between the model performance.

\begin{table}[H]
\centering
\begin{tabular}{@{}lccc@{}}
\toprule
\textbf{Template} & \textbf{Winner} & \textbf{Winner Acc.} & \textbf{Runner-up Acc.} \\
\midrule
T1 Minimal & Gemini-2.0-flash & 88.50\% & 88.00\% \\
T2 Benchmark & GPT-4o-mini & 90.00\% & 89.50\% \\
T3 Natural & Gemini-2.0-flash & 90.00\% & 89.50\% \\
T4 Order/Phrasing & GPT-4o-mini & 90.00\% & 89.00\% \\
\bottomrule
\end{tabular}
\caption{Prompt-level winner flips for the sociology extension.}
\label{tab:sociology_case}
\end{table}

\subsection{Implications for Evaluation on LLM Leaderboards}

The practical implication is that leaderboard scores would be better if accompanied by rank-stability diagnostics. A ranking based on average accuracy answers ``who scored highest under some specific setup'', while PRRR and SRRR would answer ``would the ordering survive reasonable changes to the evaluation protocol'', which is frequently a decision-relevant question as well. For benchmark reports, especially when margins are narrow, ranking robustness should be reported rather than assumed.

\section{Limitations and Risks}

This project has several limitations. First, all evaluated tasks are multiple-choice benchmarks. The findings may not transfer directly to other tasks such as generation, summarization, or open-ended reasoning. Second, we evaluate four LLMs, while the instability patterns may differ with a larger or stronger model pool. Third, our prompt perturbations are intentionally shallow and task-preserving. This design is useful for isolating non-semantic variation, but it cannot cover the full prompt design space. Fourth, our resampling protocol uses half-sized subsets without replacement; other uncertainty estimators could produce different numerical rates. Finally, API-hosted models may change over time, so exact reproducibility depends on model version stability.

\section{Conclusion}

LLM rankings could vary if we ask differently, and the pattern of ranking reversals depends on the underlying task set. ARC-Challenge remains comparatively stable, HellaSwag is strongly sensitive to prompt and subset variation, and MMLU has a stable aggregate winner but significant subject-level instability. Hence, not all leaderboards are unreliable, but leaderboard reliability must be measured. A robust evaluation report should include accuracy, uncertainty, and ranking stability metrics so that users can distinguish stable superiority from narrow, protocol-dependent ordering.

\section*{Reproducibility Statement}

\begin{itemize}[leftmargin=*]
    \item Project repository: \url{https://github.com/qianchi-zhang/LLM-Ranking-Reversals}
    %\item Model versions: \texttt{openai/gpt-4o-mini}, \texttt{google/gemini-2.0-flash-001}, \texttt{qwen/qwen-2.5-7b-instruct}, and \texttt{meta-llama/llama-3.1-8b-instruct}.
    \item Model versions used in experiments:
    \begin{enumerate}
        \item \texttt{openai/gpt-4o-mini},
        \item \texttt{google/gemini-2.0-flash-001},
        \item \texttt{qwen/qwen-2.5-7b-instruct},
        \item \texttt{meta-llama/llama-3.1-8b-instruct}.
    \end{enumerate}
    \item Main standardized pipeline: \path{run_pipeline.py} for the 900-item dataset-level experiment using \path{data/01_prompts.jsonl} and outputs under \path{data/04_analysis/}.
    \item Auxiliary expanded MMLU extension: \path{src/01_data_prep_MMLU_subject.py} prepares \path{data/01_prompts_MMLU_subjects.jsonl};    \\
    the extension reuses \path{data/02_raw_responses_MMLU_subjects.csv} and runs \path{src/04_mmlu_subject_bootstrap.py} to produce \path{data/04_mmlu_subject_analysis_expanded/}.
    \item Random seeds: main and auxiliary data sampling use seed \texttt{114514}; the expanded subject analysis uses seed \texttt{20260406}; dataset-level resampling seeds are fixed inside the Phase 4 scripts.
    %\item Main output directories used in this report: \path{data/04_analysis/}, \path{data/04_mmlu_subject_analysis_expanded/}, \path{plots/}, and \path{reports/final_report/}.
\end{itemize}

\begin{thebibliography}{9}
\bibitem{clark2018arc}
Clark, P., Cowhey, I., Etzioni, O., Khot, T., Sabharwal, A., Schoenick, C., and Tafjord, O. 2018.
\textit{Think You Have Solved Question Answering? Try ARC, the AI2 Reasoning Challenge.}

\bibitem{fourrier2024openllm}
Fourrier, C., et al. 2024.
\textit{Open LLM Leaderboard v2.} Hugging Face Blog.

\bibitem{hendrycks2020mmlu}
Hendrycks, D., Burns, C., Basart, S., Zou, A., Mazeika, M., Song, D., and Steinhardt, J. 2020.
\textit{Measuring Massive Multitask Language Understanding.}

\bibitem{miller2024errorbars}
Miller, E. 2024.
\textit{Adding Error Bars to Evals: A Statistical Approach to Language Model Evaluations.}

\bibitem{zellers2019hellaswag}
Zellers, R., Holtzman, A., Bisk, Y., Farhadi, A., and Choi, Y. 2019.
\textit{HellaSwag: Can a Machine Really Finish Your Sentence?}

\bibitem{zhu2023promptrobust}
Zhu, K., Wang, J., Zhou, J., Wang, Z., Chen, H., Wang, Y., Yang, L., Ye, W., Gong, N. Z., Zhang, Y., and Xie, X. 2023.
\textit{PromptRobust: Towards Evaluating the Robustness of Large Language Models on Adversarial Prompts.}
\end{thebibliography}

\appendix

\section{Prompt Templates}

\begin{itemize}[leftmargin=*]
    \item \texttt{minimal\_instruction}: \verb|{question}| + \verb|{choices}| + ``Select one option.''
    \item \texttt{benchmark\_style}: ``Choose the best answer to the following multiple-choice question.'' followed by question and options.
    \item \texttt{natural\_style}: ``I have a multiple-choice question for you.'' followed by question, options, and ``Which option do you think is correct?''
    \item \texttt{order\_phrasing\_variation}: ``Consider the following question.'' followed by question, options, and ``From the four options above, select the best answer.''
\end{itemize}

\iffalse
\section{Additional Output Files}

\begin{itemize}[leftmargin=*]
    \item Dataset-level PRRR: \path{data/04_analysis/prompt_ranking_reversal_rate.csv}
    \item Dataset-level SRRR: \path{data/04_analysis/subset_ranking_reversal_rate.csv}
    \item Expanded MMLU subject-level PRRR: \path{data/04_mmlu_subject_analysis_expanded/mmlu_subject_prompt_ranking_reversal_rate.csv}
    \item Expanded MMLU subject-level SRRR: \path{data/04_mmlu_subject_analysis_expanded/mmlu_subject_subset_ranking_reversal_rate.csv}
    \item Expanded MMLU subject rank probabilities: \path{data/04_mmlu_subject_analysis_expanded/mmlu_subject_rank_probability.csv}
\end{itemize}
\fi

\end{document}
